\documentclass[14pt,a4paper]{extreport}
% preamble.tex — спільна преамбула для всіх лабораторних робіт. Підключається одразу після \documentclass у кожному lab*/main.tex.
%
% ЗБІРКА: LuaLaTeX (не pdfLaTeX). Причина: пакет listings під pdfLaTeX не читає UTF-8 і зупиняє збірку на кириличних коментарях у лістингах програмного коду. LuaLaTeX обробляє UTF-8 нативно.
%
% ШРИФТ: Liberation Serif — метрично сумісний з Times New Roman і містить повний набір кириличних гліфів (TeX Gyre Termes кирилиці не має, через що український текст не відображався взагалі).

\usepackage{fontspec}
\setmainfont{Liberation Serif}
\setmonofont{DejaVu Sans Mono}[Scale=0.82]

\usepackage{polyglossia}
\setmainlanguage{ukrainian}
\setotherlanguage{english}

\usepackage{microtype}

% Поля, інтервал, відступи й заголовки нижче взяті з шаблону НУЛП, за яким написана бакалаврська робота (../bachelor-thesis/thesis/main.tex). Той самий формат приймає В.М. Шиманський і для магістерських, тож усі звіти тримаємо в ньому, а не в тому, що переказували на консультації.
\usepackage[left=25mm,right=15mm,top=15mm,bottom=23mm,footskip=1.5mm]{geometry}
\usepackage{setspace}
% Саме \setstretch{1.5}, а не \onehalfspacing: останній дає близько 1.25.
\setstretch{1.5}

\usepackage{graphicx}
\usepackage{float}
\usepackage{booktabs}
\usepackage{array}
\usepackage{tabularx}
\usepackage{longtable}
\usepackage{enumitem}
% labelindent вирівнює маркер списку з абзацним відступом — так у шаблоні НУЛП.
\setlist{nosep}
\setlist[itemize]{label=\textbullet,leftmargin=*,labelindent=1.27cm}
\setlist[enumerate]{leftmargin=*,labelindent=1.27cm}
\usepackage{indentfirst}
\usepackage{amsmath}

\usepackage{hyperref}
\hypersetup{colorlinks=true,linkcolor=black,urlcolor=blue,citecolor=black}

% xurl дозволяє \url розривати адресу в будь-якому місці, а не лише після скісної риски. Без нього довгий сегмент шляху (напр. «scopus---podrobnaya- instruktsiya-dlya-uchenogo») виходить за межі сторінки, бо дефіс для hyperref не є точкою розриву. Вантажиться після hyperref — перевизначає його \UrlBreaks.
\usepackage{xurl}

% ---------------------------- Список використаної літератури ----------------------------
% biblatex + biber зі стилем IEEE: посилання «[1]» у тексті, записи — за порядком першої згадки. Обрано IEEE, а не gost-numeric (ДСТУ 8302:2015), — так вирішив автор звітів після порівняння обох варіантів.
% Кожен звіт із джерелами сам підключає свій references.bib через \addbibresource і друкує список \printbibliography[heading=none] під власним \section. Звіт без джерел від завантаженого biblatex нічого не втрачає.
% Оточення списку — стандартне для стилю: мітка від лівого поля, текст опису вирівняний після неї. Не підганяємо його під абзацний відступ нумерованих переліків тексту (labelindent у \setlist вище) — список літератури є окремим розділом зі своїм стандартним оформленням.
\usepackage[backend=biber,style=ieee]{biblatex}
% biblatex за замовчуванням додає розтяжний проміжок після «/ . : ? =» в адресах, і з вимкненими переносами рядок вирівнюється розрідженням самої адреси («https : / / www . scopus»). Рвати адресу й так дозволяє xurl, тож проміжок лише шкодить.
\setlength{\biburlbigskip}{0mu}
% Без додаткового інтервалу між описами: розділ має йти суцільним текстом, як решта звіту. Так само в шаблоні бакалаврської.
\setlength{\bibitemsep}{0pt}
\setlength{\bibparsep}{0pt}

% ---------------------------- Лістинги програмного коду ---------------------------- Використано minted (через Pygments), а не listings: listings не вміє коректно обробляти багатобайтові символи й «висмикує» кириличні коментарі з коду, зліплюючи їх в один рядок. Потребує прапорця -shell-escape.
\usepackage{xcolor}
\usepackage[newfloat]{minted}
% labelsep=period дає «Рис. 2.1. Назва» замість «Рис. 2: Назва» — формат НУЛП.
\usepackage[labelsep=period]{caption}

% baselinestretch=1 обов'язковий: без нього лістинг успадковує півторашку основного тексту, і довгий фрагмент коду перестає вміщатися у сторінку. Усередині [H]-флоата це не помилка збірки, а нескінченний цикл порожніх сторінок — lab1 так намолотив 161 тисячу сторінок.
\setminted{
  baselinestretch=1,
  fontsize=\footnotesize,
  linenos=true,
  frame=single,
  framesep=4pt,
  breaklines=true,
  breakanywhere=true,
  tabsize=4,
  numbersep=7pt,
  bgcolor=gray!8,
}
\setminted[python]{python3=true}

\SetupFloatingEnvironment{listing}{name=Лістинг,listname=Перелік лістингів}

% ---------------------------- Нумерація та оформлення заголовків ----------------------------
\usepackage{titlesec}

\setcounter{secnumdepth}{3}
\renewcommand{\thesection}{\arabic{section}}
\renewcommand{\thesubsection}{\arabic{section}.\arabic{subsection}}
\renewcommand{\thesubsubsection}{\arabic{section}.\arabic{subsection}.\arabic{subsubsection}}

% Розділ по центру, 16 пт, великими літерами; слово «РОЗДІЛ» опущене, бо звіт не дисертація. Номер лишається: \numberwithin прив'язує до нього нумерацію рисунків («Рис. 2.1»), і без нумерованих розділів вона б зламалася. \MakeUppercase четвертим аргументом отримує текст заголовка, тому в джерелі заголовки пишуться звичайним регістром.
\titleformat{\section}[block]
  {\centering\bfseries\fontsize{16}{19.2}\selectfont}
  {\thesection.}{0.3em}{\MakeUppercase}
% Підрозділи — зліва з абзацного відступу, звичайним регістром.
\titleformat{\subsection}[block]{\bfseries}{\thesubsection.}{1em}{}
\titleformat{\subsubsection}[block]{\bfseries\itshape}{\thesubsubsection.}{1em}{}

% Незіркові \titlespacing — обов'язкові. Без них titlesec застосовує зірковий варіант, який прибирає абзацний відступ у першого абзацу після заголовка й перекриває indentfirst. Значення — типові для report.
\titlespacing{\section}{0pt}{20pt}{6pt}
\titlespacing{\subsection}{1.27cm}{6pt}{0pt}
\titlespacing{\subsubsection}{1.27cm}{6pt}{0pt}

\setlength{\parindent}{1.27cm}
\setlength{\parskip}{0pt}
\renewcommand{\arraystretch}{1.3}

% Переноси слів вимкнені — так у шаблоні НУЛП. Натомість абзац розганяється міжслівними пробілами, тому tolerance і emergencystretch підняті: без них рядки виходили б за праве поле.
\hyphenpenalty=10000
\exhyphenpenalty=10000
\tolerance=9999
\setlength{\emergencystretch}{3em}

% ---------------------------- Підписи до рисунків і таблиць ----------------------------
\captionsetup[figure]{name=Рис.,justification=centering,font={bf,it}}
\captionsetup[table]{name=Таблиця,justification=raggedleft,singlelinecheck=false,labelfont=bf,textfont=normalfont}

% Рисунки, таблиці й формули нумеруються в межах розділу: «Рис. 2.1», а не «Рис. 5». Так у шаблоні НУЛП, і так само виглядали приклади на консультації.
\numberwithin{equation}{section}
\numberwithin{figure}{section}
\numberwithin{table}{section}

% ---------------------------- Колонтитули ----------------------------
\usepackage{fancyhdr}
\pagestyle{fancy}
\fancyhf{}
\fancyfoot[R]{\thepage}
\renewcommand{\headrulewidth}{0pt}
\fancypagestyle{plain}{%
  \fancyhf{}
  \fancyfoot[R]{\thepage}
  \renewcommand{\headrulewidth}{0pt}
}

% titlepage.tex — титульний аркуш, спільний для всіх дисциплін.
% Потребує визначених у subject.tex команд:
%   \SubjectName, \LecturerName, \StudentName, \StudentGroup, \AcademicYear
%
% Використання у тексті: \MakeLabTitle{1}{Тема роботи}
%
% Рід роботи. Дисципліни з практичними роботами перевизначають його після
% % preamble.tex — спільна преамбула для всіх лабораторних робіт. Підключається одразу після \documentclass у кожному lab*/main.tex.
%
% ЗБІРКА: LuaLaTeX (не pdfLaTeX). Причина: пакет listings під pdfLaTeX не читає UTF-8 і зупиняє збірку на кириличних коментарях у лістингах програмного коду. LuaLaTeX обробляє UTF-8 нативно.
%
% ШРИФТ: Liberation Serif — метрично сумісний з Times New Roman і містить повний набір кириличних гліфів (TeX Gyre Termes кирилиці не має, через що український текст не відображався взагалі).

\usepackage{fontspec}
\setmainfont{Liberation Serif}
\setmonofont{DejaVu Sans Mono}[Scale=0.82]

\usepackage{polyglossia}
\setmainlanguage{ukrainian}
\setotherlanguage{english}

\usepackage{microtype}

% Поля, інтервал, відступи й заголовки нижче взяті з шаблону НУЛП, за яким написана бакалаврська робота (../bachelor-thesis/thesis/main.tex). Той самий формат приймає В.М. Шиманський і для магістерських, тож усі звіти тримаємо в ньому, а не в тому, що переказували на консультації.
\usepackage[left=25mm,right=15mm,top=15mm,bottom=23mm,footskip=1.5mm]{geometry}
\usepackage{setspace}
% Саме \setstretch{1.5}, а не \onehalfspacing: останній дає близько 1.25.
\setstretch{1.5}

\usepackage{graphicx}
\usepackage{float}
\usepackage{booktabs}
\usepackage{array}
\usepackage{tabularx}
\usepackage{longtable}
\usepackage{enumitem}
% labelindent вирівнює маркер списку з абзацним відступом — так у шаблоні НУЛП.
\setlist{nosep}
\setlist[itemize]{label=\textbullet,leftmargin=*,labelindent=1.27cm}
\setlist[enumerate]{leftmargin=*,labelindent=1.27cm}
\usepackage{indentfirst}
\usepackage{amsmath}

\usepackage{hyperref}
\hypersetup{colorlinks=true,linkcolor=black,urlcolor=blue,citecolor=black}

% xurl дозволяє \url розривати адресу в будь-якому місці, а не лише після скісної риски. Без нього довгий сегмент шляху (напр. «scopus---podrobnaya- instruktsiya-dlya-uchenogo») виходить за межі сторінки, бо дефіс для hyperref не є точкою розриву. Вантажиться після hyperref — перевизначає його \UrlBreaks.
\usepackage{xurl}

% ---------------------------- Список використаної літератури ----------------------------
% biblatex + biber зі стилем IEEE: посилання «[1]» у тексті, записи — за порядком першої згадки. Обрано IEEE, а не gost-numeric (ДСТУ 8302:2015), — так вирішив автор звітів після порівняння обох варіантів.
% Кожен звіт із джерелами сам підключає свій references.bib через \addbibresource і друкує список \printbibliography[heading=none] під власним \section. Звіт без джерел від завантаженого biblatex нічого не втрачає.
% Оточення списку — стандартне для стилю: мітка від лівого поля, текст опису вирівняний після неї. Не підганяємо його під абзацний відступ нумерованих переліків тексту (labelindent у \setlist вище) — список літератури є окремим розділом зі своїм стандартним оформленням.
\usepackage[backend=biber,style=ieee]{biblatex}
% biblatex за замовчуванням додає розтяжний проміжок після «/ . : ? =» в адресах, і з вимкненими переносами рядок вирівнюється розрідженням самої адреси («https : / / www . scopus»). Рвати адресу й так дозволяє xurl, тож проміжок лише шкодить.
\setlength{\biburlbigskip}{0mu}
% Без додаткового інтервалу між описами: розділ має йти суцільним текстом, як решта звіту. Так само в шаблоні бакалаврської.
\setlength{\bibitemsep}{0pt}
\setlength{\bibparsep}{0pt}

% ---------------------------- Лістинги програмного коду ---------------------------- Використано minted (через Pygments), а не listings: listings не вміє коректно обробляти багатобайтові символи й «висмикує» кириличні коментарі з коду, зліплюючи їх в один рядок. Потребує прапорця -shell-escape.
\usepackage{xcolor}
\usepackage[newfloat]{minted}
% labelsep=period дає «Рис. 2.1. Назва» замість «Рис. 2: Назва» — формат НУЛП.
\usepackage[labelsep=period]{caption}

% baselinestretch=1 обов'язковий: без нього лістинг успадковує півторашку основного тексту, і довгий фрагмент коду перестає вміщатися у сторінку. Усередині [H]-флоата це не помилка збірки, а нескінченний цикл порожніх сторінок — lab1 так намолотив 161 тисячу сторінок.
\setminted{
  baselinestretch=1,
  fontsize=\footnotesize,
  linenos=true,
  frame=single,
  framesep=4pt,
  breaklines=true,
  breakanywhere=true,
  tabsize=4,
  numbersep=7pt,
  bgcolor=gray!8,
}
\setminted[python]{python3=true}

\SetupFloatingEnvironment{listing}{name=Лістинг,listname=Перелік лістингів}

% ---------------------------- Нумерація та оформлення заголовків ----------------------------
\usepackage{titlesec}

\setcounter{secnumdepth}{3}
\renewcommand{\thesection}{\arabic{section}}
\renewcommand{\thesubsection}{\arabic{section}.\arabic{subsection}}
\renewcommand{\thesubsubsection}{\arabic{section}.\arabic{subsection}.\arabic{subsubsection}}

% Розділ по центру, 16 пт, великими літерами; слово «РОЗДІЛ» опущене, бо звіт не дисертація. Номер лишається: \numberwithin прив'язує до нього нумерацію рисунків («Рис. 2.1»), і без нумерованих розділів вона б зламалася. \MakeUppercase четвертим аргументом отримує текст заголовка, тому в джерелі заголовки пишуться звичайним регістром.
\titleformat{\section}[block]
  {\centering\bfseries\fontsize{16}{19.2}\selectfont}
  {\thesection.}{0.3em}{\MakeUppercase}
% Підрозділи — зліва з абзацного відступу, звичайним регістром.
\titleformat{\subsection}[block]{\bfseries}{\thesubsection.}{1em}{}
\titleformat{\subsubsection}[block]{\bfseries\itshape}{\thesubsubsection.}{1em}{}

% Незіркові \titlespacing — обов'язкові. Без них titlesec застосовує зірковий варіант, який прибирає абзацний відступ у першого абзацу після заголовка й перекриває indentfirst. Значення — типові для report.
\titlespacing{\section}{0pt}{20pt}{6pt}
\titlespacing{\subsection}{1.27cm}{6pt}{0pt}
\titlespacing{\subsubsection}{1.27cm}{6pt}{0pt}

\setlength{\parindent}{1.27cm}
\setlength{\parskip}{0pt}
\renewcommand{\arraystretch}{1.3}

% Переноси слів вимкнені — так у шаблоні НУЛП. Натомість абзац розганяється міжслівними пробілами, тому tolerance і emergencystretch підняті: без них рядки виходили б за праве поле.
\hyphenpenalty=10000
\exhyphenpenalty=10000
\tolerance=9999
\setlength{\emergencystretch}{3em}

% ---------------------------- Підписи до рисунків і таблиць ----------------------------
\captionsetup[figure]{name=Рис.,justification=centering,font={bf,it}}
\captionsetup[table]{name=Таблиця,justification=raggedleft,singlelinecheck=false,labelfont=bf,textfont=normalfont}

% Рисунки, таблиці й формули нумеруються в межах розділу: «Рис. 2.1», а не «Рис. 5». Так у шаблоні НУЛП, і так само виглядали приклади на консультації.
\numberwithin{equation}{section}
\numberwithin{figure}{section}
\numberwithin{table}{section}

% ---------------------------- Колонтитули ----------------------------
\usepackage{fancyhdr}
\pagestyle{fancy}
\fancyhf{}
\fancyfoot[R]{\thepage}
\renewcommand{\headrulewidth}{0pt}
\fancypagestyle{plain}{%
  \fancyhf{}
  \fancyfoot[R]{\thepage}
  \renewcommand{\headrulewidth}{0pt}
}

\input{../../../../titlepage}
: \renewcommand{\WorkKind}{ПРАКТИЧНА РОБОТА}.
\newcommand{\WorkKind}{ЛАБОРАТОРНА РОБОТА}

\newcommand{\MakeLabTitle}[2]{%
\begin{titlepage}
% Титулка — не суцільний текст, а фіксована верстка з явних \\[...] проміжків,
% тож міжрядковий інтервал основного тексту (\setstretch{1.5} у preamble.tex)
% до неї не застосовний: він розтягував кожен із ~15 рядків на 20 %, і блок
% «Виконав» разом із «Львів --- рік» переповнював сторінку. Інтервал локальний
% для оточення titlepage, тому на решту документа це не впливає.
\setstretch{1}
\begin{center}
  {\large\bfseries МІНІСТЕРСТВО ОСВІТИ І НАУКИ УКРАЇНИ}\\[6pt]
  {\large\bfseries Національний університет «Львівська політехніка»}\\[6pt]
  {\large\bfseries Інститут комп'ютерних наук та інформаційних технологій}\\[6pt]
  {\large\bfseries Кафедра «Системи штучного інтелекту»}\\[1.2cm]

  \includegraphics[width=5cm]{../../../../logo}\\[1.2cm]

  {\Large\bfseries \WorkKind{} №#1}\\[0.5cm]
  з дисципліни\\[4pt]
  «\SubjectName»\\[0.6cm]
  на тему\\[4pt]
  \textit{#2}
\end{center}

% Відступ до блоку «Виконав» гнучкий, а не фіксований (було 2 cm). Фіксований
% разом із довгим родом роботи чи темою на два рядки переповнював сторінку, і
% «Львів --- рік» зі свого \vfill з'їжджав на окрему другу сторінку. \vfill
% стискається до нуля, коли місця бракує, тож переповнення стало неможливим.
\vfill

\begin{flushright}
  {\bfseries Виконав:}\\
  студент групи \StudentGroup\\
  \StudentName\\[1cm]
  {\bfseries Викладач:}\\
  \LecturerName\\
\end{flushright}

\vfill
\begin{center}
  Львів --- \AcademicYear
\end{center}
\end{titlepage}
\clearpage
}


% Метадані дисципліни. Використовуються титульним аркушем (titlepage.tex).

\newcommand{\SubjectName}{Науковий процес та професійна науково-дослідна робота}
% Два рядки: посада з науковим ступенем, далі прізвище. `\\` розриває рядок усередині flushright-блоку титульного аркуша. Посаду (професор, не доцент) підтвердив сам викладач.
\newcommand{\LecturerName}{проф. каф. СШІ д.т.н.,\\ Ізонін І.В.}
\newcommand{\StudentName}{Галич Максим}
\newcommand{\StudentGroup}{КНСШ-13}
\newcommand{\AcademicYear}{2026}

% Рід роботи для титульного аркуша — як у practice1 (див. коментар там).
\renewcommand{\WorkKind}{ЗВІТ\\[2pt]\normalfont\large про виконання практичної роботи}
% Рисунки — скриншоти, тобто рукотворний невідтворюваний вхід, тому лежать поряд із текстом і комітяться.
\graphicspath{{assets/}}
% Два файли джерел: власні джерела звіту та бібліотека Zotero (колекція PR3). Другий перезаписує автоекспорт Better BibTeX, тому нічого рукотворного в ньому тримати не можна. Шляхи — відносно каталогу роботи, куди make render-report робить cd.
\addbibresource{references.bib}
\addbibresource{zotero.bib}

% Рисунок-скриншот на всю ширину — те саме оформлення, що й у practice2.
\newcommand{\Screenshot}[4][]{%
  \begin{figure}[H]
    \centering
    \includegraphics[width=\linewidth,#1]{#2}
    \caption{#3}
    \label{#4}
  \end{figure}
}

\begin{document}

\MakeLabTitle{3}{Пошук наукової літератури за допомогою спеціалізованих пошукових систем та бібліографічних баз}

\textbf{Мета роботи.} Навчитися здійснювати пошук наукової літератури за допомогою спеціалізованих пошукових систем та бібліографічних баз (Scopus, Google Академія), релевантної темі магістерської роботи.

\section{Постановка завдання}

Здійснити пошук за ключовими словами наукових праць, релевантних темі магістерської роботи, у пошуковій системі Google Академія та наукометричній базі Scopus~\cite{guide}. Під час пошуку в базі Scopus сформувати запит, який дає невелику кількість (до 50) релевантних темі праць. Зокрема:

\begin{enumerate}
\item записати тему та мету магістерської роботи, її об'єкт і предмет дослідження, визначити 3--6 ключових слів, які найповніше описують тему, та оформити їх у таблиці українською та англійською мовами;

\item знайти в Google Академії~\cite{scholar} щонайменше три найбільш цитовані статті за темою українською мовою, опубліковані за останній рік, виключивши з пошуку патенти та цитати, і додати їх до Zotero;

\item на основі ключових слів сформувати пошуковий запит у базі Scopus~\cite{scopus}, відфільтрувати результати за відкритим доступом, типом документа, мовою та роком публікації, навести запит у звіті, обрати щонайменше три статті та додати їх до Zotero;

\item описати проведену роботу, проілюструвати її рисунками й таблицями та навести рисунок сформованої в Zotero бібліографії у стилі IEEE;

\item зробити висновок: що зроблено, як саме і як це допоможе під час написання магістерської роботи.
\end{enumerate}

\section{Опис проведеної роботи}

\subsection{Тема магістерської роботи та ключові слова}

Тема магістерської роботи — «Удосконалення методу синтезу мовлення з керуванням голосом за текстовим описом для української мови».

\textbf{Мета роботи} — підвищити відповідність синтезованого українського мовлення текстовому опису голосу шляхом удосконалення методу синтезу мовлення з керуванням голосом за текстовим описом.

\textbf{Об'єкт дослідження} — процес синтезу мовлення з керуванням характеристиками голосу.

\textbf{Предмет дослідження} — методи та моделі синтезу мовлення з керуванням голосом за текстовим описом для української мови.

Ключові слова, які найповніше описують тему, наведено в таблиці~\ref{tab:keywords}. Англомовні відповідники взято з термінології, якою користуються автори публікацій у цій галузі, а синоніми поєднуються в запиті оператором OR.

\begin{table}[H]
\caption{Ключові слова теми магістерської роботи}
\label{tab:keywords}
\small
\begin{tabularx}{\linewidth}{|c|X|X|}
\hline
\textbf{№} & \textbf{Українською} & \textbf{Англійською} \\
\hline
1 & синтез мовлення & speech synthesis, text-to-speech (TTS) \\
\hline
2 & текстовий опис голосу & text description, natural language description, text prompt \\
\hline
3 & керування голосом & voice control, voice design, style control \\
\hline
4 & українська мова & Ukrainian \\
\hline
5 & мова з обмеженими ресурсами & low-resource language \\
\hline
\end{tabularx}
\end{table}

\subsection{Пошук у Google Академії}

Пошук у Google Академії~\cite{scholar} виконано українською мовою, щоб отримати україномовні праці за темою. Запит поступово звужувався, а після кожного кроку фіксувалася приблизна кількість результатів, яку показує система (таблиця~\ref{tab:scholar}).

\begin{table}[H]
\caption{Поетапне звуження пошуку в Google Академії}
\label{tab:scholar}
\small
\begin{tabularx}{\linewidth}{|c|X|r|}
\hline
\textbf{Крок} & \textbf{Запит і фільтри} & \textbf{Результатів} \\
\hline
1 & \texttt{синтез мовлення}, публікації з 2026~р. & 3\,130 \\
\hline
2 & фраза в лапках: \texttt{"синтез мовлення"} & 76 \\
\hline
3 & виключено патенти & 76 \\
\hline
4 & виключено цитування & 73 \\
\hline
\end{tabularx}
\end{table}

Без лапок Google Академія шукає слова запиту окремо, тому на першому кроці у видачу потрапили праці, де «синтез» стосується зовсім іншого, наприклад синтезу навчального відеоконтенту (рисунок~\ref{fig:scholar-step1}). Фраза в лапках вимагає, щоб слова стояли поруч і саме в такому порядку, і скоротила видачу в сорок разів. Фільтр за роком обмежує пошук публікаціями з 2026~р., тобто за останній рік, як вимагають методичні рекомендації.

\Screenshot{scholar-step1}{Перший крок пошуку в Google Академії: слова запиту без лапок}{fig:scholar-step1}

Наостанок з видачі виключено патенти та цитування. Цитуваннями Google Академія називає записи про праці, яких вона не знайшла в мережі, а знає лише з чужих списків літератури, тож повного тексту в них немає. Підсумкову видачу з усіма фільтрами наведено на рисунку~\ref{fig:scholar-final}.

\Screenshot{scholar-final}{Підсумкова видача Google Академії: фраза в лапках, публікації з 2026~р., без патентів і цитувань}{fig:scholar-final}

Сортувати видачу за кількістю цитувань Google Академія не дозволяє, проте впорядкування за відповідністю, яке діє за замовчуванням, цитування враховує. Обмеження останнім роком робить вимогу щодо цитованості майже формальною: праці 2026~р. ще не встигли процитувати, і лише одна з них, Гай і Шевцова~\cite{gay2026}, має одне цитування. Тому серед праць видачі обрано найцитованішу та дві найближчі до теми магістерської роботи, усі з доступним повним текстом.

Перед додаванням до бібліотеки кожну працю перевірено на тип публікації. Дві перші праці видачі розміщено в електронному архіві ХНУРЕ, і перша з них виявилася пояснювальною запискою до бакалаврської кваліфікаційної роботи, а не статтею, тому її пропущено. Відібрані праці:

\begin{itemize}
\item Іщенко та ін.~\cite{ishchenko2026} — система синтезу українського мовлення на основі моделі VITS і власного аудіокорпусу, найближча до теми магістерської роботи;
\item Білас і Різник~\cite{bilas2026} — стаття в журналі «Науковий вісник НЛТУ України» про метод генерування мовлення з передаванням емоцій, тобто з керуванням характеристиками голосу;
\item Гай і Шевцова~\cite{gay2026} — тези про інтеграцію синтезу мовлення в мобільні системи для читання, єдина праця видачі з цитуванням.
\end{itemize}

Праці додано до колекції \textit{PR3} бібліотеки Zotero через розширення Zotero Connector, яке на сторінці видачі Google Академії дозволяє вибрати потрібні записи зі списку. Метадані, які формує Google Академія, виявилися неповними й місцями хибними: усі три записи збережено як журнальні статті, а для статті Іщенко та ін. автором зазначено «Corpus» — Google Академія прочитала як ім'я автора слова з англомовної анотації. Записи виправлено вручну за титульними сторінками праць і вихідними даними збірників: двом працям зі збірників конференцій змінено тип на \textit{Conference Paper}, внесено справжніх авторів, назви збірників, місця та видавців.

\subsection{Пошук у Scopus}

Пошук у базі Scopus~\cite{scopus} виконано англійською мовою через обліковий запис Національного університету «Львівська політехніка». Scopus не містить повних текстів праць, лише їхні метадані, тому пошук ведеться в полях назви, анотації та ключових слів (\textit{Article title, Abstract, Keywords}). Як і в Google Академії, спершу налаштовувалася текстова частина запиту, а потім застосовувалися фільтри; кількість документів після кожного кроку наведено в таблиці~\ref{tab:scopus}.

\begin{table}[H]
\caption{Поетапне формування запиту в Scopus}
\label{tab:scopus}
\small
\begin{tabularx}{\linewidth}{|c|X|r|}
\hline
\textbf{Крок} & \textbf{Зміна запиту або фільтрів} & \textbf{Документів} \\
\hline
1 & синтез мовлення AND текстовий опис чи підказка (\textit{text description, natural language description, text prompt, voice design, style prompt}) & 96 \\
\hline
2 & додано синонім \textit{TTS}, форми множини та фрази \textit{voice prompt, voice description, speaker description, natural language instruction} & 108 \\
\hline
3 & прибрано \textit{TTS}, \textit{voice prompt} і \textit{natural language description}; \textit{natural language} — лише поруч зі словами \textit{voice, speaker, style} & 103 \\
\hline
4 & тип документа — стаття, мова — англійська, відкритий доступ & 9 \\
\hline
5 & роки публікації 2024--2026 & 7 \\
\hline
6 & тип документа — стаття або доповідь на конференції & 29 \\
\hline
\end{tabularx}
\end{table}

Перший варіант запиту дав 96 документів, проте більшість із них стосувалася не синтезу мовлення з керуванням голосом (рисунок~\ref{fig:scopus-step1}). Фраза \textit{text description} трапляється в будь-якій праці, де щось описують текстом, тож у видачу потрапили, наприклад, системи генерування підписів до зображень.

\Screenshot{scopus-step1}{Перший варіант запиту в Scopus}{fig:scopus-step1}

Спроба розширити запит синонімами на другому кроці лише додала шуму. Скорочення \textit{TTS} і фраза \textit{voice prompt} притаманні допоміжним системам для людей з порушеннями зору, які озвучують розпізнаний текст чи знак і подають користувачеві голосові підказки, а \textit{natural language description} так само позначає й текстовий опис зображення. На третьому кроці ці терміни прибрано, а фразу \textit{natural language} залишено лише за умови, що не далі ніж за три слова від неї стоїть \textit{voice}, \textit{speaker} або \textit{style} — для цього Scopus має оператор близькості \texttt{W/n}. Після цього серед перших результатів з'явилися праці саме за темою (рисунок~\ref{fig:scopus-step3}).

\Screenshot{scopus-step3}{Уточнений запит в Scopus з оператором близькості}{fig:scopus-step3}

Далі застосовано фільтри, рекомендовані методичними рекомендаціями: лише статті, лише англійською мовою, лише з відкритим доступом і лише за останні три роки. Залишилося сім документів, і жоден з них не стосувався теми: це праці про голосових асистентів та озвучення навчальних матеріалів (рисунок~\ref{fig:scopus-articles}). Причина в тому, що праці із синтезу мовлення публікуються переважно в матеріалах конференцій (ICASSP, Interspeech, ACL, AAAI): серед 103 документів третього кроку 79 були доповідями на конференціях. Тому до типу документа, крім статей, додано доповіді на конференціях; решту фільтрів збережено.

\Screenshot{scopus-articles-only}{Видача Scopus за фільтром «лише статті»: релевантних праць немає}{fig:scopus-articles}

Підсумковий запит дав 29 документів, що відповідає вимозі отримати до 50 релевантних праць. Scopus формує його в розширеному синтаксисі разом з усіма фільтрами (рисунок~\ref{fig:scopus-query}):

\begin{minted}[breaklines,fontsize=\small]{text}
TITLE-ABS-KEY ( ( "speech synthesis" OR "text-to-speech" ) AND ( "text prompt*" OR "style prompt*" OR "voice description*" OR "speaker description*" OR "natural language instruction*" OR "voice design" OR ( "natural language" W/3 ( voice OR speaker OR style ) ) ) ) AND ( LIMIT-TO ( LANGUAGE , "English" ) ) AND ( LIMIT-TO ( OA , "all" ) ) AND ( LIMIT-TO ( DOCTYPE , "cp" ) OR LIMIT-TO ( DOCTYPE , "ar" ) ) AND PUBYEAR > 2023 AND PUBYEAR < 2027
\end{minted}

Запит відтворюваний: будь-хто, виконавши його, отримає ту саму видачу на дату пошуку. Цим пошук у наукометричній базі відрізняється від пошуку в загальній пошуковій системі чи запиту до великої мовної моделі, і саме тому такий запит можна навести в розділі магістерської роботи з оглядом літератури як частину методики відбору джерел.

\Screenshot{scopus-query}{Підсумковий запит Scopus у розширеному синтаксисі та перші документи видачі}{fig:scopus-query}

Фільтр відкритого доступу на рисунку~\ref{fig:scopus-query} не видно, оскільки панель фільтрів довша за екран; його наведено окремо на рисунку~\ref{fig:scopus-oa} разом із частиною видачі.

\Screenshot{scopus-open-access}{Фільтр відкритого доступу та продовження підсумкової видачі Scopus}{fig:scopus-oa}

З підсумкової видачі обрано три доповіді, найближчі до теми магістерської роботи:

\begin{itemize}
\item Zhou та ін.~\cite{zhou2026} — IndexTTS2, модель синтезу мовлення з незалежним керуванням тембром, емоціями та тривалістю, у якій емоцію можна задати текстовим описом;
\item Qiang та ін.~\cite{qiang2026} — UniSonate, єдина модель для генерування мовлення, музики та звукових ефектів за текстовими інструкціями;
\item Li та ін.~\cite{li2026} — DisCo-Speech, керований синтез мовлення для голосу, не баченого під час навчання, на основі кодека з розділеними характеристиками мовлення.
\end{itemize}

Праці додано до колекції \textit{PR3} через Zotero Connector. Метадані зі Scopus виявилися значно якіснішими, ніж з Google Академії: тип запису, повний перелік авторів, сторінки та DOI заповнено правильно. Виправлено лише назву збірника, яку Scopus записує скорочено («Proc. Annu. Meet. Assoc. Comput Linguist.»), — її замінено повною назвою з поля назви конференції.

\subsection{Бібліографія у стилі IEEE}

Бібліографію до шести праць колекції \textit{PR3} сформовано в Zotero у стилі IEEE~\cite{ieee-guide} командою \textit{Create Bibliography from Items}, експортовано у формат RTF і відкрито в текстовому процесорі LibreOffice Writer (рисунок~\ref{fig:ieee}). Стиль нумерує записи в порядку першої згадки, наводить ініціали перед прізвищем, назву праці бере в лапки, а назву збірника чи журналу виділяє курсивом. Цей самий стиль використовує й список літератури цього звіту, який формується з тієї ж бібліотеки через автоматичний експорт Better BibTeX у файл \texttt{zotero.bib}, як і в попередній практичній роботі.

\Screenshot{zotero-ieee}{Бібліографія колекції \textit{PR3}, сформована в Zotero у стилі IEEE, у LibreOffice Writer}{fig:ieee}

\section{Результати}

У результаті роботи знайдено шість праць, релевантних темі магістерської роботи: три україномовні праці в Google Академії та три англомовні доповіді в базі Scopus. Усі праці додано до колекції \textit{PR3} бібліотеки Zotero; їх зведено в таблиці~\ref{tab:results}.

\begin{table}[H]
\caption{Праці, знайдені за темою магістерської роботи}
\label{tab:results}
\small
\begin{tabularx}{\linewidth}{|c|l|X|l|c|}
\hline
\textbf{№} & \textbf{Джерело} & \textbf{Праця} & \textbf{Тип} & \textbf{Цитувань} \\
\hline
1 & Google Академія & Іщенко та ін., VITS~\cite{ishchenko2026} & доповідь & 0 \\
\hline
2 & Google Академія & Білас, Різник~\cite{bilas2026} & стаття & 0 \\
\hline
3 & Google Академія & Гай, Шевцова~\cite{gay2026} & тези & 1 \\
\hline
4 & Scopus & Zhou та ін., IndexTTS2~\cite{zhou2026} & доповідь & 1 \\
\hline
5 & Scopus & Qiang та ін., UniSonate~\cite{qiang2026} & доповідь & 0 \\
\hline
6 & Scopus & Li та ін., DisCo-Speech~\cite{li2026} & доповідь & 0 \\
\hline
\end{tabularx}
\end{table}

Підсумки пошуку в обох системах порівняно в таблиці~\ref{tab:compare}.

\begin{table}[H]
\caption{Порівняння пошуку в Google Академії та Scopus}
\label{tab:compare}
\small
\begin{tabularx}{\linewidth}{|l|X|X|}
\hline
\textbf{Ознака} & \textbf{Google Академія} & \textbf{Scopus} \\
\hline
Мова запиту & українська & англійська \\
\hline
Початкова видача & 3\,130 & 96 \\
\hline
Підсумкова видача & 73 & 29 \\
\hline
Засоби звуження & фраза в лапках, рік, виключення патентів і цитувань & оператори AND/OR, оператор близькості W/n, фільтри за роком, типом документа, мовою, відкритим доступом \\
\hline
Відтворюваність & запит не експортується & запит експортується в розширеному синтаксисі \\
\hline
Якість метаданих & неповні, потребували ручного виправлення & повні, з DOI \\
\hline
Покриття & зокрема кваліфікаційні роботи й тези, які треба відсіювати & лише праці з індексованих видань \\
\hline
\end{tabularx}
\end{table}

\section{Висновки}

У ході практичної роботи сформульовано тему, мету, об'єкт і предмет магістерської роботи, визначено ключові слова теми українською та англійською мовами, виконано пошук за ними в Google Академії та наукометричній базі Scopus і додано по три знайдені праці з кожної системи до бібліотеки Zotero, з якої сформовано бібліографію у стилі IEEE.

\begin{enumerate}[ref=\arabic*]
\item Ефективний пошук будується поступовим звуженням запиту. Точна фраза в лапках скоротила видачу Google Академії з 3\,130 до 76 праць за один крок, а в Scopus оператор близькості відсіяв праці, у яких потрібні слова трапляються, але в іншому значенні. Розширення запиту синонімами не завжди корисне: синонім, поширений в іншій галузі, лише додає шуму, тому кожну зміну запиту варто перевіряти за першими результатами видачі.

\item Фільтри, рекомендовані методичними рекомендаціями, потрібно узгоджувати з особливостями галузі. Обмеження лише статтями в журналах залишило в Scopus сім праць, жодна з яких не стосувалася теми, оскільки дослідження із синтезу мовлення публікуються переважно в матеріалах конференцій. Додавання доповідей на конференціях дало 29 релевантних документів.

\item Google Академія індексує значно ширше коло матеріалів, ніж Scopus, зокрема кваліфікаційні роботи студентів і тези конференцій, тому тип кожної знайденої праці треба перевіряти перед використанням. Її метадані також менш надійні: у двох із трьох записів довелося виправляти тип, а в одному — авторів. Метадані Scopus, навпаки, повні й містять DOI.

\item Вимога обрати найбільш цитовані праці за останній рік суперечлива: праці, опубліковані цього року, ще не встигли набрати цитувань, і максимум у видачі обох систем становив одне цитування. Для огляду літератури в магістерській роботі доцільно розширити період пошуку до трьох років, звичайних для галузі штучного інтелекту, і окремо шукати найцитованіші праці за весь час, на які спирається галузь.

\item Пошук англійською мовою в Scopus дав праці, значно ближчі до теми магістерської роботи, ніж україномовний пошук: україномовних праць саме про керування голосом за текстовим описом за останній рік не знайдено, а знайдені праці присвячені синтезу українського мовлення загалом. Це підтверджує, що тема має наукову новизну для української мови, а найближчі аналоги слід шукати в англомовних виданнях.

\item Для магістерської роботи найціннішим результатом є відтворюваний запит Scopus. Його можна навести в розділі з оглядом літератури як частину методики відбору джерел, розширити на ширший період і повторювати в міру появи нових публікацій, а знайдені праці, як-от IndexTTS2 з керуванням емоцією за текстовим описом, стають відправною точкою для огляду методів-аналогів.
\end{enumerate}

\section{Список використаної літератури}

\printbibliography[heading=none]

\end{document}
